\subsection{Examples of duality}

PROPOSITION 17. --- Let \(n \geqslant 1\) be an integer. Denote by \(\boldsymbol{\mu}_{n}\) the group of \(n\) -th roots of unity in \(\mathbf{C}\) . The groups \(\mathbf{Z}/n\mathbf{Z}\) and \(\boldsymbol{\mu}_{n}\) are in duality with respect to the map induced by passage to the quotient of the map \(\mathbf{Z} \times \boldsymbol{\mu}_{n} \to \mathbf{U}\) defined by \((m,z) \mapsto z^{m}\) .

The group \(\mathbf{Z}/n\mathbf{Z}\) coincides with the set of homomorphisms \(\chi\) from \(\mathbf{Z}/n\mathbf{Z}\) to \(\mathbf{U}\) . These are of the form \(m \mapsto \chi(1)^m\) where \(\chi(1)\) is an arbitrary element of \(\mathbf{U}\) such that \(\chi(1)^n = 1\) , whence the result.

COROLLARY 1. --- Let G be a finite commutative group. The dual group \(\widehat{\mathbf{G}}\) is isomorphic to G.

The group G is isomorphic to a finite product of cyclic groups (A, VII, p.~22, th. 3), and its dual group is isomorphic to the product of the dual groups of these (prop. 2 of II, p.~206). One is therefore reduced to the case where G is cyclic, which follows from proposition 17 since the group \(\mu_{n}\) is cyclic of order n (A, V, p.~75, th. 1).

COROLLARY 2. --- Let G be a locally compact commutative group. The group G is finite if and only if \(\widehat{\mathbf{G}}\) is finite. A closed subgroup H of G has finite index if and only if its orthogonal is finite.

By duality, the first assertion follows from the fact that the dual of a finite group is finite (corollary 1). The second follows from it, since \(\widehat{\mathrm{G / H}}\) is identified with \(\mathrm{H}^{\perp}\) (th. 4 of II, p.~226).

PROPOSITION 18. --- For \(\widehat{\mathbf{G}}\) to be compact, it is necessary and sufficient that \(\mathbf{G}\) be discrete. If \(\mathbf{G}\) is discrete, the dual measure of the counting measure on \(\mathbf{G}\) is the normalized Haar measure on \(\widehat{\mathbf{G}}\) . If \(\mathbf{G}\) is compact, the dual measure of the normalized Haar measure is the counting measure on \(\widehat{\mathbf{G}}\) .

Suppose G is discrete, and let \(\alpha\) be the counting measure on G. Let \(\varphi\) be the characteristic function of \(e \in G\) . We have \(\mathcal{F}_{\mathrm{G}}(\varphi) = 1\) on \(\widehat{G}\) . Since \(\mathcal{F}_{\mathrm{G}}(\varphi)\) tends to 0 at infinity, the group \(\widehat{G}\) is compact.

Furthermore, for the dual measure \(\widehat{\alpha}\) of \(\alpha\), the function \(\mathcal{F}_{\mathrm{G}}(\varphi)\) must have integral \(\varphi(e)=1\) (prop. 12 of II, p.~219). Therefore \(\widehat{\alpha}(\widehat{\mathrm{G}})=1\) .

Suppose G is compact. Then the Haar measure dx belongs to \(\mathcal{M}^{1}(G)\) . Its Fourier transform is strictly positive at \(\chi = e\) and zero for \(\chi \neq 0\) (prop. 6 of II, p.~210). Since it is continuous on \(\widehat{G}\) , the group \(\widehat{G}\) is discrete. If the measure of G is 1, we deduce by duality from the previous case that the dual measure of the measure dx is the counting measure on \(\widehat{G}\) .

COROLLARY 1 (Orthogonality relations). --- Suppose G is discrete and equipped with the counting measure. For x and y in G, we have

\[
\int_ {\widehat {G}} \chi (x) \overline {{\chi (y)}} d \chi = \left\{ \begin{array}{l l} 0 & \text{if } x \neq y \\ 1 & \text{if } x = y. \end{array} \right.
\]

This follows from the cor. of theo. 1 of II, p.~215 and from duality.

COROLLARY 2. --- Let H be a closed subgroup of G.

\begin{enumerate}
\def\labelenumi{\alph{enumi})}
\item
  For H to be compact, it is necessary and sufficient that \(H^{\perp}\) be open in \(\widehat{G}\) ;
\item
  For H to be open, it is necessary and sufficient that \(H^{\perp}\) be compact in \(\widehat{G}\) .
\item
  To say that \(H^{\perp}\) is open amounts to saying that \(\widehat{G}/H^{\perp}\) is discrete, but \(\widehat{G}/H^{\perp}\) is isomorphic to \(\widehat{H}\) (th. 4 of II, p.~226); the assertion therefore follows from prop. 18. Assertion b) follows by duality from assertion a) applied to \(H^{\perp}\) .
\end{enumerate}

COROLLARY 3. --- Let \((\mathrm{H}_{i})_{i\in\mathrm{I}}\) be a decreasing filtered family of compact subgroups of G. For G to be identified with the projective limit of the groups \(G/H_{i}\) , it is necessary and sufficient that \(\hat{G}\) be the union of the open subgroups \(H_{i}^{\perp}\) .

To say that G is identified with the projective limit of the \(G/H_{i}\) amounts to saying that \(\bigcap_{i}H_{i}=\{e\}\) (TG, III, p.~60, prop. 2), that is, that \(\bigcup_{i}H_{i}^{\perp}\) is dense in \(\widehat{G}\) (cor. 4 of th. 4 of II, p.~226). Now \(\bigcup_{i}H_{i}^{\perp}\) is an open, hence closed, subgroup of \(\widehat{G}\) .

COROLLARY 4. --- Let I be a set and let \((\mathrm{H}_{i})_{i\in\mathrm{I}}\) be a family of compact groups. The dual of the product group of \(H_{i}\) is the discrete group direct sum of the groups \(\widehat{H}_{i}\) .

This is a special case of Cor. 3.

PROPOSITION 19. --- Let K be a locally compact, non-discrete field, not necessarily commutative, and let G be the additive group of K, whose group law is written additively. Let \(\chi\) be a unitary character of G distinct from 1. For \(x,y\in\mathbf{G}\), set \(\theta(x,y)=\chi(xy)\). Then G is in duality with itself with respect to \(\theta\) .

For \(y \in \mathbf{G}\), let \(\chi_y\) be the map from \(\mathbf{G}\) into \(\mathbf{U}\) such that \(\chi_y(x) = \chi(xy)\). We have \(\chi_y \in \widehat{\mathbf{G}}\), and it must be shown that \(\beta: y \mapsto \chi_y\) is a topological group isomorphism from \(\mathbf{G}\) into \(\widehat{\mathbf{G}}\).

The map \(\beta\) is an injective homomorphism from G into \(\widehat{\mathbf{G}}\); it is continuous (TG, X, p.~28, th. 3 applied to the continuous map \(\theta\) of \(\mathbf{G} \times \mathbf{G}\) into \(\mathbf{C}\)). Let us prove that \(\theta\) is a homeomorphism onto its image. It suffices (Lemma 2 of II, p.~200) to show that for every neighborhood U of 0 in K, there exists a neighborhood V of \(e\) in \(\widehat{\mathbf{G}}\) such that \(\overline{\beta}^{-1}(\mathrm{V}) \subset \mathrm{U}\). Let \(x \mapsto |x|\) be an absolute value on K defining the topology of K (AC, VI, §9, no. 1, prop. 1), and let \(x_0 \in \mathrm{K}\) be such that \(\chi(x_{0}) \neq 1\); let us denote \(\eta = |\chi(x_{0}) - 1| > 0\). Let U be a neighborhood of 0 in K. There exists \(\delta > 0\) such that U contains the set of \(y \in K\) such that \(|y| < \delta\). Let V be the set of characters \(\xi \in \widehat{G}\) such that \(|\langle\xi, x\rangle - 1| < \eta\) for every element \(x \in K\) satisfying \(|x| \leqslant |x_{0}|/\delta\). It is a neighborhood of e in \(\widehat{G}\). If \(y \neq 0\) is such that \(\beta(y)\) belongs to V, we therefore have \(|\chi(xy) - 1| < |\chi(x_{0}) - 1|\) for all x such that \(|x| \leqslant |x_{0}|/\delta\). Consequently, we have \(|x_{0}y^{-1}| > |x_{0}|/\delta\), and hence \(|y| < \delta\), so that \(y \in U\).

Since \(\beta\) is a homeomorphism onto its image, its image is closed in \(\widehat{\mathbf{G}}\) (TG, III, p.~22, cor. 2). But on the other hand, the orthogonal of the image of \(\beta\) is the set of elements \(x\) of \(\mathbf{G}\) such that \(\chi(xy) = 1\) for all \(y \in \mathbf{G}\), and is therefore reduced to \(\{0\}\). The image of \(\beta\) is therefore dense in \(\widehat{\mathbf{G}}\) (corollary 3 of II, p.~228). We conclude that \(\beta\) is surjective.

COROLLARY 1. --- Let K be a locally compact, non-discrete, not necessarily commutative field, and \(\chi\) a nontrivial unitary character of the additive group of K. Let E be a finite-dimensional topological vector space over K. The map \(\theta\) from \(E \times E'\) into U defined by \(\theta(x,\lambda)=\chi(\langle\lambda,x\rangle)\) for \((\lambda,x)\in\mathrm{E}'\times\mathrm{E}\) puts the topological groups E and E' in duality.

Let n be the dimension of E and \((e_{1},\ldots,e_{n})\) a basis of E. It allows one to identify E and \(K^{n}\) (EVT, I, p.~14, th. 2). The result then follows from prop. 19 and prop. 2 of II, p.~206.

We denote by T the group R/Z.

COROLLARY 2. --- a) The group R is in duality with itself with respect to the mapping \((x,y)\mapsto\exp(2i\pi xy)\) , and the dual measure of Lebesgue measure is Lebesgue measure;

\begin{enumerate}
\def\labelenumi{\alph{enumi})}
\setcounter{enumi}{1}
\tightlist
\item
  The groups Z and T are in duality with respect to the map obtained by passing to the quotient from the map from \(Z \times R\) to U such that \((n, x) \mapsto \exp(2i\pi nx)\) . The dual Haar measure of the counting measure on Z is the normalized Haar measure on R/Z.
\end{enumerate}

The group \(\mathbf{R}\) is in duality with itself relative to the mapping \((x,y)\mapsto \exp (2i\pi xy)\) by prop. 19. Let us identify \(\widehat{\mathbf{R}}\) with \(\mathbf{R}\) . The orthogonal of \(\mathbf{Z}\) in \(\widehat{\mathbf{R}} = \mathbf{R}\) is then \(\mathbf{Z}\) , and \(b)\) follows from th. 4 of II, p.~226.

Let \(\alpha\) be the counting measure on Z and \(\gamma\) the normalized Haar measure on T. If \(\beta\) denotes Lebesgue measure on \(\mathbf{R}\), we have \(\gamma = \beta/\alpha\) since these two Haar measures on R/Z have mass 1. The Haar measure \(\widehat{\alpha}\) on \(\widehat{Z} = T\) is the normalized Haar measure (prop. 18), and the Haar measure \(\widehat{\gamma}\) is the counting measure on \(\mathbf{Z}\) (loc. cit.). By Prop. 16 of II, p.~231, the dual measure of \(\beta\) is therefore the measure \(\beta\) .

Remark. --- In particular, one recovers the determination of \(\mathsf{X}(\mathsf{L}^{1}(\mathbf{Z}))\) as done in Example 4 of I, p.~36.

For every integer \(n \geqslant 0\) and \((x, y) \in \mathbf{R}^n \times \mathbf{R}^n\) , we denote by

\[
x \cdot y = \sum_ {j = 1} ^ {n} x _ {j} y _ {j}.
\]

COROLLARY 3. --- Let \(n \geqslant 1\) an integer. The group \(\mathbf{R}^{n}\) is in duality with itself relative to the mapping \((x,y) \mapsto \exp(2i\pi x \cdot y)\) and the dual measure of the Lebesgue measure on \(\mathbf{R}^{n}\) is the Lebesgue measure. The groups \(\mathbf{Z}^{n}\) and \(\mathbf{T}^{n} = \mathbf{R}^{n}/\mathbf{Z}^{n}\) are in duality with respect to the map obtained by passing to the quotient from the map \((n,x) \mapsto \exp(2i\pi x \cdot y)\) , and the dual Haar measure of the counting measure on \(\mathbf{Z}^{n}\) is the normalized Haar measure on \((\mathbf{R}/\mathbf{Z})^{n}\) .

This follows from Lemma 7 of II, p.~225, from Proposition 10 of II, p.~217, and from Corollary 2.

Remark. --- Given a subgroup H of \(R^{n}\) , there corresponds its orthogonal \(H^{\perp}\) , a subgroup of \(\widehat{R^{n}} = R^{n}\) , which is none other than the subgroup associated with H defined in TG, VII, p.~6, no. 3.

In what follows, we shall identify the dual of \(R^{n}\) (resp. of \(T^{n}\) ) with \(R^{n}\) (resp. with \(Z^{n}\) ) by means of the duality of the corollary. In particular, for \(f \in \mathrm{L}^{1}(\mathbf{R}^{n})\) , its Fourier transform is identified with the function on \(R^{n}\) with values in C which, to \(y \in R^{n}\) , associates

\[
\mathcal {F} (f) (y) = \int_ {\mathbf {R} ^ {n}} f (x) \exp (- 2 i \pi x \cdot y) d x.
\]

COROLLARY 4. --- The group \(\mathbf{R}^{*}\) is in duality with the group \(\{-1,1\} \times \mathbf{R}\) by the map \((x,(\sigma,t)) \mapsto \sigma(x/|x|)|x|^{it}\) . The group \(\mathbf{R}_{+}^{*}\) is in duality with \(\mathbf{R}\) by the map \((x,t) \mapsto x^{it}\) .

Indeed, the map \(x \mapsto (x/|x|, \log(|x|))\) is a topological group isomorphism from \(R^{*}\) onto \(\{-1, 1\} \times R\) . The assertion then follows from lemma 7 of II, p.~225, from corollary 2, and from the fact that the unitary characters of \(\{-1, 1\}\) are 1 and \(x \mapsto x\) .

Let \(p\) be a prime number. The field \(\mathbf{Q}_{p}\) of \(p\)-adic numbers is the completion of \(\mathbf{Q}\) with respect to the \(p\)-adic valuation (INT, VII, § 1, no. 6, example, and AC, VI, § 3, no. 4, example 4). For every \(x \in \mathbf{Q}_{p}\), there exists a unique integer \(\nu\geqslant0\) and a unique integer q satisfying \(0\leqslant q<p^{\nu}\) such that \(qp^{-\nu}-x\in Z_{p}\) (A, VII, p.~10, th. 2, applied to the principal ideal domain \(Z_{p}\) and to the set \(R_{p}\) of integers j such that \(0\leqslant j<p\)). We denote by \(\lambda(x)=qp^{-\nu}\) .

PROPOSITION 20. — The mapping \(x \mapsto \exp(2i\pi\lambda(x))\) is a unitary character of \(\mathbf{Q}_{p}\) whose kernel is \(\mathbf{Z}_{p}\) .

For \(x_1\) and \(x_2\) in \(\mathbf{Q}_p\), by definition \(\lambda(x_1 + x_2) - \lambda(x_1) - \lambda(x_2) \in \mathbf{Z}_p \cap \mathbf{Q} = \mathbf{Z}\). The map \(\lambda\) is moreover locally constant since \(\lambda(x + y) = \lambda(x)\) if \(y \in \mathbf{Z}_p\). It then follows that \(x \mapsto \exp(2i\pi\lambda(x))\) is a unitary character of \(\mathbf{Q}_p\) . Since \(\lambda(x) \in \mathbf{Z}\) if and only if \(x \in \mathbf{Z}_p\) , the kernel of this character is \(\mathbf{Z}_p\) .

Recall that one calls normalized Haar measure on the additive group of \(Q_{p}\) the unique Haar measure \(\mu\) such that \(\mu(\mathbf{Z}_{p}) = 1\) (INT, VII, §1, no. 6, example).

COROLLARY. --- a) The group \(\mathbf{Q}_{p}\) is in duality with itself relative to the mapping \((x,y)\mapsto\exp(2i\pi\lambda(xy))\). The normalized Haar measure on \(\mathbf{Q}_{p}\) is then its own dual;

\begin{enumerate}
\def\labelenumi{\alph{enumi})}
\setcounter{enumi}{1}
\tightlist
\item
  The groups \(Z_{p}\) and \(Q_{p}/Z_{p}\) are in duality with respect to the map obtained by passing to the quotient from the map defined by \((z,x)\mapsto\exp(2i\pi\lambda(zx))\) , and the dual measure of the normalized Haar measure on \(Z_{p}\) is the counting measure on \(Q_{p}/Z_{p}\) .
\end{enumerate}

The proof follows step by step that of Cor. 2 of Prop. 19.

